\chapter{Tartışma}
\label{ch:tartisma}

Elde edilen deneysel bulgular, kademeli tıbbi teşhis sistemlerinin klinik entegrasyon potansiyelini açıkça göstermektedir.

\section{Kademeli Yaklaşımın Klinik Faydaları}
Tek aşamalı derin öğrenme modelleri tüm görüntüleri en karmaşık modelden geçirerek yüksek bir enerji ve zaman maliyeti oluşturmaktadır. Önerilen kademeli model ise vakaların \%65'ten fazlasını hafif olan MobileNetV2 modeli ile çözebilmekte, böylece sunucu maliyetlerinde ve karbon emisyonunda ciddi bir azalma sağlamaktadır. 

En kritik durumlarda ise (belirsizliğin yüksek olduğu vakalar), model otomatik olarak kararı daha güçlü olan Tier 2 modeline yönlendirmektedir. Bu sayede hem işlem gücü korunmakta hem de klinik güvenilirlik maksimum düzeye çıkarılmaktadır.

\section{Sınırlılıklar ve Gelecek Çalışmalar}
Sistemin temel sınırlılığı, yönlendirme oranının seçilen güven eşiğine ($\tau$) oldukça duyarlı olmasıdır. Gelecek çalışmalarda, dinamik eşik belirleyen pekiştirmeli öğrenme ajanlarının entegrasyonu araştırılacaktır. Ayrıca, çoklu etiketli ve diğer akciğer patolojilerini de kapsayacak genişletilmiş kademeli yapılar üzerinde çalışılması planlanmaktadır.
